\chapter{Layer 7: Policy Reasoning (Answer-Set Programming)}
\label{ch:policy}

\section{Purpose and Contract}
\label{sec:policy-purpose}

Layer~7 answers the question \emph{Is this result permitted?} Its inbound
contract is ``a set of facts and a policy''; its outbound contract is ``a
stable model (yes/no + justification) or \texttt{inconsistent}.'' It is
implemented with Answer-Set Programming (ASP) via Clingo.

\section{Why Non-Monotonic Reasoning}
\label{sec:policy-asp}

Mathematical truth is monotonic: adding axioms never retracts theorems.
\emph{Policy}, by contrast, is non-monotonic: adding a new fact (``this agent
is revoked'') can retract a previously valid permission. ASP models exactly
this: a program's answer sets are its stable models, and additional
constraints can eliminate models. The \texttt{layers/asp/} module therefore
hosts the governance decisions the coherence model of \cref{ch:coherence}
sketches at the repository level.

\section{The Module Contract}
\label{sec:policy-module}

The \texttt{README} for the layer states:
\begin{quote}
\textbf{Layer 7: ASP Stable Models} \\
Purpose: Answer Set Programming (Clingo) for non-monotonic reasoning. \\
Status: Stable.
\end{quote}
``Stable'' here means the solver converges to a consistent set of answer sets
for the current policy corpus. A result that has passed L1--L6 is only
released if at least one stable model permits it.

\section{Example Policy Shape}
\label{sec:policy-example}

%\examplepolicy

A representative policy fragment (illustrative) couples an artifact to its
provenance and to an agent's clearance:
\begin{verbatim}
% Facts: artifact A was produced by agent X with receipt R
produced(a1, agent_kitty, r1).
cleared(agent_kitty, publish).

% Rule: an artifact may publish iff its producer is cleared
% and its receipt is sealed
may_publish(A) :- produced(A, X, R),
                 cleared(X, publish),
                 sealed(R).

% Constraint: never publish an unsealed artifact
:- produced(A, X, R), not sealed(R), published(A).
\end{verbatim}
Clingo computes the stable models; if \texttt{may\_publish(a1)} holds in some
model and no constraint is violated, the artifact advances. Otherwise it is
held at the boundary.

\section{Connection to Governance}
\label{sec:policy-gov}

Layer~7 is where the abstract governance kernel ($\ensuremath{Omega}\ensuremath{leftarrow}
\bigodot\ensuremath{wedge}\boxdot\ensuremath{wedge}\ensuremath{dots}$) becomes executable. The entropy and
intercoil predicates of \cref{ch:coherence} are, conceptually, ASP facts; the
``no component advances unless coherent'' rule is an ASP constraint. The
pipeline thus inherits the constellation's governance at the granularity of a
single artifact.

\section{Policy as Answer-Set Constraints}
\label{sec:policy-asp}

Concretely, a policy such as ``no unverified numeric claim may be
published'' is rendered as the ASP rule set:
\begin{verbatim}
verified(X) :- proof(X, lean).
published(X) :- claim(X).
:- published(X), not verified(X).
\end{verbatim}
The solver reports UNSAT for any schedule that publishes an unverified
claim, and SAT (with the published set) otherwise. This keeps policy
enforcement inside the same P-time verifier family as the P/NP layer
(\cref{ch:pnp}).
